%---------------------------------------------------------------------------%
%->> Backmatter
%---------------------------------------------------------------------------%
\chapter[致谢]{致\quad 谢}\chaptermark{致\quad 谢}% syntax: \chapter[目录]{标题}\chaptermark{页眉}
%\thispagestyle{noheaderstyle}% 如果需要移除当前页的页眉
%\pagestyle{noheaderstyle}% 如果需要移除整章的页眉

三年的硕士学习即将结束，论文写到这里，也意味着这一阶段的学习与研究工作接近尾声。回顾这段时间，从课题的确定、平台适配工作的推进，到实验反复调试和论文最终成稿，过程中得到了许多老师、同学和家人的帮助与支持。在此，谨向所有关心和帮助过我的人致以诚挚的谢意。

首先，我要衷心感谢我的导师张福新老师。论文从选题确定到后续研究推进，张老师都给予了我耐心而细致的指导。无论是技术路线的把握、实验方案的设计，还是论文内容的修改完善，张老师都提出了许多宝贵意见。张老师严谨务实的治学态度和认真负责的工作作风，让我在硕士阶段受益匪浅，也让我对科研工作有了更加踏实的认识。

感谢中国科学院计算技术研究所为我提供良好的学习和科研环境。感谢课题组的老师和同学们，特别是朱奇正师兄，在课题推进过程中给予我的帮助。无论是在 ART 运行时与 LoongArch 平台适配方面，还是在实验环境搭建、问题排查和结果分析过程中，大家都给了我很多具体而实际的支持。这些交流和讨论让我收获很多，也帮助我顺利完成了论文工作

同时，感谢课题相关前期工作的积累，以及开源社区提供的软件基础、技术资料和文档支持，使得本文的研究工作能够在已有基础上继续展开。对在平台搭建、工具维护和资料整理过程中提供帮助的老师和同学，在此一并表示感谢。

最后，我要感谢我的家人。感谢他们一直以来的理解、支持与陪伴。在学习和论文写作过程中，正是家人的鼓励与包容，使我能够更加专注地完成硕士阶段的研究工作。

谨以此文，向所有帮助过我的人致谢。

\chapter{作者简历及攻读学位期间发表的学术论文与研究成果}

\section*{作者简历：}

2019 年 09 月——2023 年 07 月,在中国科学院大学计算机学院获得学士学位。
2023 年 09 月——2026 年 07 月,在中国科学院大学计算机学院攻读硕士学位。

\section*{已发表（或正式接受）的学术论文：}

\section*{申请或已获得的专利：}

\section*{参加的研究项目及获奖情况：}

{
\setlist[enumerate]{}% restore default behavior
\begin{enumerate}[nosep]
    \item 参与基于 LoongArch 架构的 Android 运行时适配与执行引擎优化相关研究，完成 AOSP 15 ART
    执行引擎适配、验证链路建立以及基准测试驱动的性能分析与优化工作。
\end{enumerate}
}

\cleardoublepage[plain]% 让文档总是结束于偶数页，可根据需要设定页眉页脚样式，如 [noheaderstyle]
%---------------------------------------------------------------------------%